% AI-produced by GPT-6 Astra (OpenAI, ChatGPT Pro mode). No priority claim.
\subsection{Ordering a sufficiently long repaired necklace}

We first specify the set-colimit assertion used in the ordering argument.
The target in this subsection is exactly the covered repaired target:
it has the prescribed fixed points, private points, and actual-length
path targets. No additional unused points are required.

\begin{lemma}[Proved: the repaired target is the strict necklace colimit]
\label{C:necklace-set-colimit}
Let $f,g:S\rightrightarrows T$ be injections of finite sets. Suppose
the nonidentity part of the relation $f(x)\mapsto g(x)$ has only directed
paths. Let $n\geq2$ bound their numbers of vertices, choose $k\geq n$,
and put $s=4k+1$. Let
\[
 \alpha_i:T\longrightarrow Z_k,\qquad i\in\mathbb Z/s,
\]
be the actual-path-length repaired beta maps. Define
\[
 (\lambda_i,\mu_i)=
 \begin{cases}
 (g,f),&0\leq i<2k,\\
 (g,g),&i=2k,\\
 (f,g),&2k<i<4k,\\
 (f,f),&i=4k.
 \end{cases}
\]
Then $Z_k$ is the coequalizer in sets of
\[
 \coprod_{i\in\mathbb Z/s}S
 \rightrightarrows
 \coprod_{i\in\mathbb Z/s}T,
 \qquad
 (i,x)\longmapsto(i,\lambda_i x),\quad
 (i,x)\longmapsto(i+1,\mu_i x).
 \tag{NC}\label{eq:akr-necklace-coequalizer}
\]
In particular, compatible maps from the displayed copies of $T$ to
any set induce a unique map from $Z_k$.
\end{lemma}
\begin{proof}
Write a nontrivial relation path as
$t_0\mapsto t_1\mapsto\cdots\mapsto t_{m-1}$, where $m\geq2$.
Its target consists of the $s$ points $0,\ldots,4k$, and the restriction
of copy $i$ is
\[
 \beta_i^m(t)=
 \begin{cases}
 t+i,&i\leq2k,\ t+i\leq2k,\\
 t+i+2k-m+1,&i\leq2k,\ t+i>2k,\\
 t-i+6k-m+2,&i>2k,\ t-i+6k-m+2>2k,\\
 t-i+4k,&i>2k,\ t-i+6k-m+2\leq2k.
 \end{cases}
\]
Substitution gives $\alpha_i\lambda_i=\alpha_{i+1}\mu_i$.
It also gives the exact consecutive overlap: for $i\leq2k$ it is
the image under $\beta_i^m$ of $\{1,\ldots,m-1\}$, and for $i>2k$
it is the image of $\{0,\ldots,m-2\}$. These are exactly the images
of the appropriate source edges of the path under $\lambda_i$.

For clarity, every identification in the target is generated by these
neighboring identifications. For a target point $u$, its visibility
set $I_u=\{i:u\in\beta_i^m(\{0,\ldots,m-1\})\}$ is
\[
\begin{array}{c|c}
\text{range of }u&I_u\\\hline
0\leq u\leq m-2&[4k-u,4k]\cup[0,u]\\
m-1\leq u\leq2k&[u-m+1,u]\\
2k+1\leq u\leq4k-m+1&[6k-m+2-u,6k+1-u]\\
4k-m+2\leq u\leq4k&[u-2k,6k+1-u].
\end{array}
\]
Each row is a nonempty cyclic interval. Because every $\beta_i^m$
is injective, there is one preimage of $u$ at each index of $I_u$.
The consecutive-overlap calculation joins exactly the adjacent
preimages. Thus every such fiber is connected by the relations in
\eqref{eq:akr-necklace-coequalizer}.

A fixed point has one preimage in every copy, joined cyclically by
those same relations. A private point has exactly one preimage.
The three kinds of targets are disjoint, and the visibility table
also proves that all path-target points are covered. Therefore
$\coprod_iT\to Z_k$ is onto, and each of its fibers is exactly one
equivalence class generated by the two displayed maps. This is the
stated strict coequalizer.
\end{proof}

\begin{proposition}[Proved: all sufficiently long repaired necklaces are orderable]
\label{C:ordered-necklace}
Let $S,T$ be finite linearly ordered sets and let
$f,g:S\rightrightarrows T$ be increasing injections. Fix $n\geq2$
at least the number of vertices in every nontrivial path of the relation
$f(x)\mapsto g(x)$. There is an integer $K\geq n$ such that, for every
$k\geq K$, the covered repaired target $Z_k$ admits a linear order
making every generating map $\alpha_i:T\to Z_k$ increasing.
The prescribed order on $T$ is unchanged. The assertion permits any
additional lower bounds on $k$.
\end{proposition}
\begin{proof}
The partial bijection $f(S)\to g(S)$, $f(x)\mapsto g(x)$, is
increasing. It has no nontrivial directed cycle: restriction to such
a cycle would give a nonidentity increasing permutation of a finite
linear order. Its nonfixed components are therefore paths, so
Lemma~\ref{C:necklace-set-colimit} applies.

Embed $T$ increasingly in $\mathbb R$. Extend the increasing partial
bijection to an increasing homeomorphism
\[
 h:\mathbb R\longrightarrow\mathbb R,
 \qquad h(f(x))=g(x)\quad(x\in S).
 \tag{H}\label{eq:akr-order-h}
\]
For nonempty $S$, this can be done by linear interpolation through
the finitely many points $(f(x),g(x))$, with slope-one translations
on the two unbounded complementary intervals. If $S$ is empty, take
$h$ to be the identity.

Every component $J$ of $\mathbb R\setminus\operatorname{Fix}(h)$
is invariant under $h$. On $J$, either $h(t)>t$ everywhere or
$h(t)<t$ everywhere. In the first case $h^r(t)$ tends to the right
endpoint of $J$ as $r\to\infty$; in the second it tends to the left
endpoint. Indeed the iterates are monotone, and a finite limit in the
interior would be fixed by continuity, contrary to the definition of
$J$. The assertion includes convergence to an infinite endpoint.

For each $x\in S$, the points $f(x)$ and $g(x)$ either coincide in
$\operatorname{Fix}(h)$ or lie in the same such component $J$.
In the nonfixed case put
\[
 B_J=(f(S)\cup g(S))\cap J.
\]
Only finitely many components meet these finite sets. Consequently
there is $K\geq n$ such that, for every $k\geq K$,
\[
 \begin{cases}
 h^{2k}(g(S)\cap J)>\max B_J,&h|_J>\operatorname{id},\\
 h^{2k}(g(S)\cap J)<\min B_J,&h|_J<\operatorname{id}.
 \end{cases}
 \tag{S}\label{eq:akr-order-separation}
\]
An inequality for a finite set means that every element satisfies it.
Once these inequalities hold, monotonicity of the iterates makes them
hold for every larger $k$.

Fix such a $k$. There is an increasing homeomorphism $w$ of
$\mathbb R$ satisfying simultaneously
\[
 w(g(x))=f(x)\quad(x\in S),\qquad
 w(h^{2k}g(x))=h^{2k}g(x)\quad(x\in S).
 \tag{W}\label{eq:akr-order-w}
\]
To verify this, first consider the prescription on its finite domain.
Within a positive component $J$, all points of $h^{2k}g(S)\cap J$
lie strictly above both $f(S)\cap J$ and $g(S)\cap J$; in a negative
component they lie strictly below both. Thus the increasing bijection
$g(x)\mapsto f(x)$ is compatible with fixing these additional points.
Prescriptions in different components preserve their relative order.
At fixed points the two prescriptions agree, since $g(x)$ fixed implies
$f(x)=g(x)$ by injectivity of $h$. The finite prescription is therefore
well-defined and strictly increasing. Linear interpolation, with
slope-one tails, extends it to the required homeomorphism $w$.

Define increasing maps $a_i:T\to\mathbb R$ by
\[
 a_i=
 \begin{cases}
 h^i|_T,&0\leq i\leq2k,\\
 (w\circ h^{4k+1-i})|_T,&2k<i\leq4k.
 \end{cases}
\]
They satisfy exactly the prescribed necklace equations. For
$0\leq i<2k$, equation~\eqref{eq:akr-order-h} gives
\[
 a_i g=h^i g=h^{i+1}f=a_{i+1}f.
\]
At the upper junction, equation~\eqref{eq:akr-order-w} gives
\[
 a_{2k}g=h^{2k}g=w h^{2k}g=a_{2k+1}g.
\]
For $2k<i<4k$,
\[
 a_i f=w h^{4k+1-i}f=w h^{4k-i}g=a_{i+1}g.
\]
At the lower junction,
\[
 a_{4k}f=w h f=w g=f=a_0f.
\]
All equalities are equalities of actual functions on $S$.

The strict set-colimit property of
Lemma~\ref{C:necklace-set-colimit} now gives a unique function
\[
 \rho:Z_k\longrightarrow\mathbb R,
 \qquad \rho\alpha_i=a_i.
\]
Order the finite set $Z_k$ first by its $\rho$-values and then by any
chosen linear order within each fiber of $\rho$. If $t<t'$ in $T$,
then $a_i(t)<a_i(t')$, so
$\rho\alpha_i(t)<\rho\alpha_i(t')$. Hence this order makes every
$\alpha_i$ strictly increasing. Possible equal values belonging to
different copies cause no problem, since no two points of one copy
have equal $\rho$-values. This proves the proposition.
\end{proof}

\begin{remark}[Proved scope and effective parameter selection]
\label{C:ordered-necklace-scope}
The proposition asserts existence for all sufficiently large $k$;
it does not assert orderability at every minimally admissible value.
Its proof is a strict set-colimit argument and uses no conversion
of a homotopy-coherent cocone into a strict one.

The choice can be made effectively. Place $T$ at integer coordinates
and choose the piecewise-linear $h$ just described, with rational
coefficients. Its fixed intervals and the endpoints of the other
components are determined by finitely many affine equations. Iterate
the finitely many rational points $g(x)$ and increase $k$ until
\eqref{eq:akr-order-separation} holds, together with all additional
numerical lower bounds. The convergence argument proves termination.
The subsequent prescriptions for $w$ and $\rho$ are finite and rational.
Thus arbitrary tie-breaking in the last step may also be replaced by
a fixed deterministic ordering rule.
\end{remark}
